% LSPV-1 REPORT::CYCLOROOTFW1::2026-10-14T01:07:55+09:00
% !TEX program = xelatex
\documentclass[11pt]{article}
\usepackage[T1]{fontenc}
\usepackage{lmodern,amsmath,amssymb,amsthm,booktabs,tabularx,xurl,hyperref}
\usepackage[margin=0.86in]{geometry}
\newtheorem{theorem}{Theorem}
\newtheorem{corollary}[theorem]{Corollary}
\newcommand{\F}{\mathbb F_2}
\newcommand{\EXACT}{\textbf{PROVED/EXACT}}
\newcommand{\COMP}{\textbf{COMPUTATIONAL}}
\newcommand{\COND}{\textbf{CONDITIONAL}}
\newcommand{\CAND}{\textbf{CANDIDATE}}
\newcommand{\OPEN}{\textbf{OPEN}}
\newcommand{\RET}{\textbf{FALSIFIED/RETIRED}}
\hypersetup{colorlinks=true,urlcolor=blue,linkcolor=blue}
\setlength{\emergencystretch}{3em}
\title{CYCLOROOTFW1: A Cyclotomic Cube Selector in the\\Superlogarithmic-Valuation Regime}
\author{UNIVERSAL\_CORE\_LOWER\_BOUND / LSPV-1}
\date{14 October 2026, 01:07 KST}
\begin{document}
\maketitle

\begin{center}
\fbox{\parbox{0.95\linewidth}{
\textbf{Executive status: P versus NP remains OPEN.}
For the unbounded sequence $d_a=2\cdot3^a$, the canonical irreducible
modulus $X^{2\cdot3^a}+X^{3^a}+1$ exposes the residue of $X$ as an
element of exact order $3^{a+1}$.  Adleman--Manders--Miller correction
therefore selects cube roots from the target alone in deterministic
polynomial time, although
$3^{\nu_3(d_a)}=d_a/2\gg\log d_a$.  The usual charged Tseitin lift then
reconstructs a full edge witness.  This upper-bound firewall is \CAND{}
pending independent review; it is not an unrestricted lower bound.}}
\end{center}

\section{Claim card}
\begin{tabularx}{\linewidth}{p{0.27\linewidth}X}
\toprule
Claim ID & CYCLOROOTFW1 \\
Status before run & Fixed valuation selector \EXACT; logarithmic-subfield extension \CAND; complementary degrees \OPEN \\
Claim attacked & Superlogarithmic $3$-primary order or unbounded $\nu_3(d)$ suffices to obstruct uniform target-only cube-root selection \\
Quantifiers & Every $a\ge0$, canonical typed degree $d_a=2\cdot3^a$, every target $b\in\mathbb F_{2^{d_a}}$, and every even connected Tseitin charge with a $d_a$-edge cotree \\
Model & One uniform deterministic random-access bit program, full input access, polynomial workspace \\
Charged resource & Recognition, field arithmetic, image test, correction, reconstruction, verification, and output \\
SAME\_REALIZATION & Same modulus, target, selected root, graph, charge, cotree, and reconstructed witness \\
Falsifier & A legal $a$ for which the trinomial is reducible, $X$ has wrong order, or a promised cube has no correction root \\
\bottomrule
\end{tabularx}

The selected obstacles were deterministic nonresidue acquisition and the
missing whole-machine transfer from nonlinear structure to time.  This run
resolves the first on one superlogarithmic infinite subtype and leaves the
second open.

\section{Primary literature and transfer classification}
Adleman--Manders--Miller introduced finite-field root extraction using a
nonresidue \cite{AMM}; Cao--Sha--Fan give detailed correction formulas
\cite{Cao}.  Sze studies deterministic root extraction without supplied
nonresidues under explicit arithmetic hypotheses \cite{Sze}.
Bhargava--Ivanyos--Mittal--Saxena connect nonresidue construction,
irreducibility, and root finding while retaining general derandomization
questions \cite{BIMS}.

AMM correction is a \textbf{directly transferable theorem mechanism}: the
residual is literally a cube equation in the represented field.  The
cyclotomic representation below is proved directly.  Generic deterministic
root finding and cryptographic root hardness are \textbf{failed transfers}:
the former is parameter-sensitive and the latter is neither unconditional
nor hard on this representation.  The cross-domain coding lesson is that
algebraic structure can expose an algorithmic weakness.

\section{Cyclotomic representation theorem}
For $a\ge0$ put
\[
 n_a=3^{a+1},\qquad d_a=\varphi(n_a)=2\cdot3^a,\qquad
 f_a(X)=X^{d_a}+X^{d_a/2}+1=\Phi_{n_a}(X).
\]

\begin{theorem}[canonical cyclotomic cube selector]
For every $a\ge0$, $f_a$ is irreducible over $\F$.  In
\[
 K_a=\F[X]/(f_a),\qquad \alpha=X\bmod f_a,
\]
$\alpha$ has exact multiplicative order $3^{a+1}$ and is a cubic
nonresidue.  One uniform deterministic algorithm, given the canonical type
$(a,d_a,f_a)$ and $b\in K_a$, returns $x$ with $x^3=b$ iff $b$ is a cube,
and otherwise rejects, in polynomial bit time and workspace.
\end{theorem}

\begin{proof}
If $k$ is odd then $2^k\not\equiv1\pmod3$.  If $k$ is even, LTE gives
\[
 \nu_3(2^k-1)=\nu_3(4^{k/2}-1)=1+\nu_3(k).
\]
Thus $3^{a+1}\mid2^k-1$ implies $2\cdot3^a\mid k$, and conversely
$2^{2\cdot3^a}\equiv1\pmod{3^{a+1}}$.  Hence
\[
 \operatorname{ord}_{3^{a+1}}(2)=2\cdot3^a=d_a. \tag{1}
\]
For a primitive $3^{a+1}$-st root $\zeta$, the Frobenius orbit length is
the least $k$ with $\zeta^{2^k}=\zeta$, namely $d_a$.  Its minimal
polynomial over $\F$ therefore has degree $d_a$, divides
$\Phi_{3^{a+1}}=f_a$, and equals $f_a$.  This proves irreducibility and
that $\alpha$ has exact order $3^{a+1}$.

Let $N=2^{d_a}-1$.  Again by LTE,
\[
 \nu_3(N)=a+1,\qquad N=3^{a+1}m,\qquad 3\nmid m. \tag{2}
\]
The exact order of $\alpha$ makes it a noncube.  Return zero for $b=0$.
For nonzero $b$, test $b^{N/3}=1$, exactly membership in the cube-image
subgroup, and reject if it fails.  Otherwise put
\[
 e=3^{-1}\bmod m,\qquad y=b^e.
\]
For cube $b$, the element $y^3/b$ lies in the subgroup generated by
$\alpha^3$, of order $3^a$.  Exactly one candidate
\[
 x_j=y\alpha^j,\qquad 0\le j<3^a=d_a/2, \tag{3}
\]
satisfies $x_j^3=b$.  Test and return it.

The recognizer recomputes $d_a$ and the three coefficients of $f_a$.
Polynomial-basis multiplication costs polynomial time in $d_a$;
exponents have $O(d_a)$ bits; and (3) has $d_a/2$ candidates.  All parsing,
image testing, correction, verification, workspace, and output costs are
polynomial in explicit input length.  No randomization, source preimage,
factorization oracle, or nonuniform advice is used.
\end{proof}

The proof is complete but new, so status is \CAND{} pending independent
review.  Since $3^{\nu_3(d_a)}=d_a/2$, this sequence lies outside every
fixed logarithmic-subfield window for large $a$.  Failure of the earlier
enumeration algorithm is therefore not a hardness theorem.

\section{Tseitin residual corollary}
Let $X_a$ be an explicitly encoded connected graph with a declared spanning
tree and an ordered cotree $C_a$ of size $d_a$, for example
$K_{2,d_a+1}$.  Let $c$ be an even charge and identify
$\pi_{C_a}x$ with $K_a$.  Define
\[
 G_{a,b,c}(x)=[B_{X_a}x=c]\land[(\pi_{C_a}x)^3=b]. \tag{4}
\]
\begin{corollary}[charged target-only section]
The typed family (4) has a uniform deterministic polynomial-time decision
and witness algorithm.  It rejects malformed records, odd charges, and
off-image targets; otherwise it selects a root by the theorem, places its
bits on the cotree, and uniquely solves the tree leaf-to-root.
\end{corollary}
\begin{proof}
Even total charge is the consistency condition for a connected incidence
system.  Once cotree values are fixed, subtracting their boundary leaves a
tree system with a unique reverse-leaf-order solution.  Recognition,
selection, reconstruction, and verification are polynomial.
\end{proof}

This does not solve an arbitrary source conjunct.  It shows only that
unbounded valuation, superlogarithmic correction order, full projected
rank, nonlinearity, proper cube image, and three-to-one nonzero fibers do
not themselves force residual work.

\section{Bounded reproducible audit}
The dependency-free program
\path{experiments/review/cyclotomic_cube_selector_firewall.py} ran once in
5.221 seconds with seed 141020260107 and passed 1,210 assertions with zero
failures:
\begin{itemize}
\item $a=0,1,2,3,4$ and degrees $2,6,18,54,162$;
\item direct polynomial-basis irreducibility, exact-order, and noncube checks;
\item 68 exhaustive target truth tables at degrees two and six;
\item 219 selected cubes, 236 off-image rejections, and 219 complete
Tseitin witness reconstructions;
\item 25 malformed-type or odd-charge rejections; and
\item correction width through $3^4=81$.
\end{itemize}
The audit is \COMP{}.  Degrees 18, 54, and 162 use sampled targets, and
finite checks do not prove the theorem.  Raw output is
\path{results/raw/review/2026-10-14_0107_cyclotomic_cube_selector.json}.
The code SHA-256 is
\texttt{d1e437be2897be1ec9286845ea610b77563595b7eb9a61036da236a3582bee15}.
No chunking was needed; timing is not extrapolated.

\section{Barrier, hostile-control, wrapper, and PQ audit}
The selector allows full target access, polynomial workspace, compact field
arithmetic, and charges recognition, exponentiation, candidate testing,
tree reconstruction, verification, and output.  DEAD/COPY/DELAY wrappers do
not remove the canonical modulus; INTERP/RECOMP can execute the selector.
Opaque affine coordinate changes are not covered unless a polynomial basis
map is declared or recovered.  Horn, Krom, affine, bounded-treewidth,
product, padding, and redundant-clause controls remain easy.

This upper bound relativizes and algebrizes; stronger semantic or algebraic
access only helps.  Natural Proofs is not invoked.  DNNF, PC/PCR,
Resolution, pebbling, one-tape, communication, and circuit lower bounds do
not transfer to unrestricted time here.  No ETH, GRH, cryptographic,
random-oracle, or average-case MQ premise is used.

For the post-quantum ledger, the cyclotomic field and public
$3$-primary-order element are an \textbf{EXACT OBJECT/THEOREM TRANSFER}.
Schur-square, quadratic-relation, Hilbert/syzygy, collision-entropy,
LWE-style extraction, and generic binding language are \textbf{FAILED
TRANSFERS} to an unrestricted lower bound here: the structure is exploitable.

\section{CORPUS REVIEWED}
This run surveyed the current GitHub tree and read root \path{main.tex},
\path{AGENTS.md}, \path{CODEX.md}, the complete PNP research skill and its
mode/protocol references, and all seven task specifications.  It read
AFFINE-RESIDUAL-REDIRECT1, EARLY22-INTEGRATION1, SINGLE22-SEPARATION1,
newest MAIN VALSUBREVIEW1, and relevant CUBEPERMFW1/PROMISEGATE1 snapshots.
Actual specialist evidence included PLANTEDCUBEFW1, CUBEROOTSEL1,
VALSUBCUBEFW1, and CHECKABLEPERM1 with cited code/results.  HEC/VSC-HEC,
G-PPC-HEC, CRS-PPC, RCCN, CCSPC, and Search-G were inspected as
restricted-model transfer guards.  EARLY22/SINGLE22 exact-$(2,2)$ compilers
and proof projections remain preserved only within their stated systems.

Predecessor recovery was unnecessary: LSPV-5's canonical
\path{research/main/2026-10-14_0048_MAIN_VALSUBREVIEW1.tex} is substantive,
has the required MAIN header and sentinel, and read back from GitHub at blob
\texttt{3ae51bbf30e11ebac5856a1cb00ab2e12f3b38ec}.

\section{DUPLICATION / CONTRADICTION AUDIT}
Repository searches for cyclotomic cube roots, $\Phi_{3^{a+1}}$, the
trinomial above, and valuation selectors found no prior CYCLOROOTFW1 theorem.
CUBEROOTSEL1 covers $a=0$; VALSUBCUBEFW1 covers each fixed $a$ by
constant-size subfield enumeration; VALSUBREVIEW1 permits growing $a$ only
when $3^a=O(\log d)$.  This uniform sequence instead has $3^a=d/2$ and
exposes the nonresidue in the representation.

There is no contradiction with general root extraction, MQ hardness, SAT
NP-completeness, or exact-$(2,2)$ restricted proof-system bounds.  This is a
tractable typed residual subtype, not a solver for its possible source
conjunct.

\section{WHAT CHANGED RELATIVE TO ROOT MAIN}
\begin{itemize}
\item The logarithmic-subfield condition is not a hardness boundary: a
canonical family far outside it has a polynomial target-only selector.
\item Unbounded $\nu_3(d)$ and $3^{\nu_3(d)}\gg\log d$ are \RET{} as
stand-alone residual-hardness premises.
\item Representation must be charged both ways: a canonical modulus can
expose the algebraic witness a lower-bound proposal hoped was hard to find.
\item The growing-factor premise and P-versus-NP status remain \OPEN{}.
\end{itemize}

\section{ADOPTED / REJECTED / RETIRED PRIOR COMPONENTS}
\textbf{Adopted:} affine elimination, typed recognition,
polynomial-basis arithmetic, cube-image verification, AMM correction, and
forest lifting.  \textbf{Preserved:} valid EARLY22/SINGLE22 compiler and
restricted proof-system results, and the uninstantiated growing-factor lemma.
\textbf{Rejected:} uncharged advice, randomized expected time, cryptographic
cube-root hardness, and failure of one algorithm as time-hardness evidence.
\textbf{Retired:} unbounded valuation or superlogarithmic $3$-primary size as
a stand-alone unavoidable-work premise.

\section{STALE OR OVERSTATED CLAIMS FOUND}
Root MAIN correctly says the superlogarithmic range escapes the previous
enumeration argument, but that must not be read as a hardness frontier.
CYCLOROOTFW1 gives a different representation-driven algorithm on a sequence
with $3^{\nu_3(d)}=d/2$.  No claim is made for arbitrary supplied irreducible
moduli at those degrees.

\section{Status and handoff}
\begin{center}
\begin{tabularx}{0.98\linewidth}{p{0.40\linewidth}p{0.19\linewidth}X}
\toprule
Claim & Status & Scope \\
\midrule
Cyclotomic irreducibility/order proof & \CAND & Complete; review pending \\
Uniform selector on $d=2\cdot3^a$ & \CAND & Canonical trinomial only \\
Finite selector/Tseitin audit & \COMP & 1,210 checks, zero failures \\
Unbounded valuation alone implies hardness & \RET & Explicit subtype refutes sufficiency \\
Growing-factor theorem & \COND & Premise uninstantiated \\
Unrestricted all-program lower bound & \OPEN & Missing implication \\
P versus NP separation & \OPEN & Load-bearing arrow remains open \\
\bottomrule
\end{tabularx}
\end{center}

\textbf{Single highest-value next target.}
Independently review the cyclotomic proof, then test whether arbitrary typed
degree-$2\cdot3^a$ moduli admit a uniformly computable isomorphism to the
canonical trinomial field.  A positive result would retire the whole
complementary-degree cube route; failure of known isomorphism algorithms
would still not be a lower bound.

\begin{thebibliography}{9}
\bibitem{AMM}
L. M. Adleman, K. L. Manders, and G. L. Miller,
\emph{On Taking Roots in Finite Fields}, FOCS 1977, 175--178.
doi:10.1109/SFCS.1977.18.
\bibitem{Cao}
Z. Cao, Q. Sha, and X. Fan,
\emph{Adleman--Manders--Miller Root Extraction Method Revisited},
arXiv:1111.4877, 2011.
\bibitem{Sze}
T.-W. Sze,
\emph{On Taking $r$-th Roots without $r$-th Nonresidues over Finite Fields
and Its Applications}, arXiv:1105.5852, 2011.
\bibitem{BIMS}
V. Bhargava, G. Ivanyos, R. Mittal, and N. Saxena,
\emph{Irreducibility and $r$-th Root Finding over Finite Fields},
arXiv:1702.00558, 2017.
\end{thebibliography}

\end{document}
% END_OF_VERIFIED_REPORT::LSPV-1::CYCLOROOTFW1

